% Regression: p3_sliding — imported current-behavior fixture from the handoff corpus.
% Formula: Sliding windows: a virtual operation on the bond e = <(3,2),(3,3)> can be
% p3_sliding — arXiv:1804.04964 (normal PEPS FT), Corollary 8 sketch, p. 24.
% Sliding windows: a virtual operation on the bond e = <(3,2),(3,3)> can be
% read as a physical operation on any of four 2x2 windows (K = L = 2), each
% obtained from the previous by sliding one lattice step (right-to-left,
% then down); the paper joins consecutive windows by arrows.
%
% Window sequence (rows x cols of the 5x5 clip):
%   W1 = (2-3, 3-4)  ->  W2 = (2-3, 2-3)  ->  W3 = (2-3, 1-2)
%   ->  W4 = (3-4, 1-2),
% the distinguished edge e fixed in all four panels.
%
% Hue mapping noted: paper draws the window red and the edge fat green;
% tenkz has slot=selected (Action red) for the window and ONE distinguished
% -edge hue (also Action red), so the paper's two-hue window/edge contrast
% collapses unless a theme rebinds a colour.
\documentclass[varwidth=30cm, border=6pt]{standalone}
\usepackage{amsmath, amssymb}
\usepackage{tenkz}
\begin{document}
\begin{tenkzcd}[column sep=large]
  \begin{tenkzlattice}[rows=5, cols=5, boundary legs, compact]
    \tnregion[slot=selected, outline]{(2-3,3-4)}
    \tnedge{(3,2)-(3,3)}
  \end{tenkzlattice}
  \arrow[r]
  &
  \begin{tenkzlattice}[rows=5, cols=5, boundary legs, compact]
    \tnregion[slot=selected, outline]{(2-3,2-3)}
    \tnedge{(3,2)-(3,3)}
  \end{tenkzlattice}
  \arrow[r]
  &
  \begin{tenkzlattice}[rows=5, cols=5, boundary legs, compact]
    \tnregion[slot=selected, outline]{(2-3,1-2)}
    \tnedge{(3,2)-(3,3)}
  \end{tenkzlattice}
  \arrow[r]
  &
  \begin{tenkzlattice}[rows=5, cols=5, boundary legs, compact]
    \tnregion[slot=selected, outline]{(3-4,1-2)}
    \tnedge{(3,2)-(3,3)}
  \end{tenkzlattice}
\end{tenkzcd}
\end{document}
